\chapter{Эквивариантная граница и невозможность топологического выбора сектора}

Предыдущая глава закрыла spin-cover bridge для уже выбранного
вращательно-эквивариантного проекторного ежа. Теперь проверяется последний
вопрос этой топологической ветви: требует ли глобальный фон проекта
существования хотя бы одного такого ежа, или однородный нулевой сектор
остаётся допустимым состоянием.

Проверка обязана быть некруговой. Нельзя начинать с проколотого пространства,
сферы вокруг дефекта или граничной карты ненулевой степени: все эти объекты
уже предполагают наличие дефекта.

\section{Явный однородный конкурент}

Текущая конечная геометрия использует постоянные матричные углы и
коэффициентные проекторы. Поэтому на любом выбранном внешнем носителе $M$
существует постоянное проекторное поле
\begin{equation}
 P_0(x)=
 \begin{pmatrix}
 1&0&0\\0&0&0\\0&0&0
 \end{pmatrix},
 \qquad x\in M.
 \label{eq:v5-sector-selection-constant-projector}
\end{equation}
Оно удовлетворяет
\begin{equation}
 P_0^2=P_0,
 \qquad
 dP_0=0,
 \qquad
 [P_0(x)]=\text{const}\in\mathbb{RP}^2.
\end{equation}
Следовательно, его проекторная градиентная энергия, локальный заряд на любой
малой сфере и индуцированный хопфов класс равны нулю.

Полная вращательная инвариантность действия не исключает этот конкурент.
Эквивариантность поля $P(Rx)=RP(x)R^{\mathsf T}$ была условием
сферического ежа, тогда как вакуумное поле может быть пространственным
скаляром с внутренней постоянной осью или усреднённым состоянием. Требовать
от одного фиксированного $P_0$ ковариантного вращения вместе с координатой
означало бы уже выбрать дефектное представление действия группы.

\begin{nogocriterion}[Симметрия действия не выбирает дефектное решение]
Из $SO(3)$-инвариантности функционала не следует, что все его стационарные
точки обязаны удовлетворять ежовой эквивариантности. Постоянный сектор
остаётся допустимым, пока его не исключают граничное условие, аномалия или
отрицательная мода полного гессиана.
\end{nogocriterion}

\section{Глобальный носитель не содержит свободного сферического заряда}

Для исторического носителя
\begin{equation}
 K=\mathbb{RP}^3\times S^1
\end{equation}
имеем
\begin{equation}
 \pi_2(K)=0,
 \qquad
 H^2(K;\mathbb Z)\simeq\mathbb Z_2.
 \label{eq:v5-sector-selection-carrier-groups}
\end{equation}
Таким образом, глобальный фон содержит только торсионный класс степени два,
но не свободный целочисленный генератор локальной сферы.

Пусть $i:S^2\hookrightarrow K$ --- малая сфера вокруг предполагаемой
точки. Для торсионного класса $a\in H^2(K;\mathbb Z)$ выполнено $2a=0$.
После ограничения
\begin{equation}
 2i^*a=i^*(2a)=0.
\end{equation}
Но $H^2(S^2;\mathbb Z)=\mathbb Z$ не имеет кручения, поэтому
\begin{equation}
 i^*a=0.
 \label{eq:v5-sector-selection-torsion-restriction}
\end{equation}
Следовательно, глобальная $\mathbb Z_2$-линия на $K$ не может сама
индуцировать локальный хопфов класс $c_1=\pm1$ на окружающей сфере.

Этот вывод согласуется со старым red-team нулевого промта: даже после
предварительного выбора $\mathbb{RP}^3$ и окружности оставались тривиальный
и нетривиальный торсионные классы, а уникальный устойчивый носитель не был
выведен.

Для позднего кандидата $S^4$ ситуация ещё жёстче:
\begin{equation}
 H^2(S^4;\mathbb Z)=0.
\end{equation}
Ни один из двух исторических глобальных фонов не заставляет локальную
сферу нести свободный класс Черна один.

\section{Глобальная классификация не создаёт дефект}

Топология нематических текстур на общем домене с заданным дефектным
множеством классифицируется скрученными когомологиями и данными накрытия.
Она различает уже существующие глобальные классы, но не утверждает, что
дефектное множество обязано быть ненулевым. См.
\path{arXiv:1604.03429}, \path{arXiv:1107.1169} и
\path{arXiv:2308.05730}.

В ветви без дисклокационных линий, где проекторное поле имеет глобальный
ориентированный подъём, заряды точечных дефектов на замкнутом
трёхмерном пространстве удовлетворяют условию нейтральности. Если из
$M$ удалить малые шары $B_i$ вокруг точек, то в
$H_2(M\setminus\cup_i B_i)$ ориентированная сумма граничных сфер равна
нулю. Спаривание с замкнутым классом поднятой карты даёт
\begin{equation}
 \sum_i Q_i=0.
 \label{eq:v5-sector-selection-charge-neutrality}
\end{equation}
Поэтому замкнутый фон скорее требует компенсации ежа анти-ежом, чем
принуждает одиночный класс $+1$. При наличии линий классификация становится
скрученной и сложнее, но тем более не возникает теоремы о неизбежном
одиночном заряде.

\section{Сверка с доказанными функционалами}

Топологическая допустимость нулевого сектора могла бы быть преодолена
динамической неустойчивостью. Однако уже построенные функционалы дают
обратный результат.

Для действия внутренней петли
\begin{equation}
 \mathcal S_{\rm loop}(V)=
 \tau\!\left([N,V]^*[N,V]\right)
\end{equation}
тождественная петля $V_0=I$ имеет
\begin{equation}
 \mathcal S_{\rm loop}(V_0)=0,
\end{equation}
тогда как минимум внутри сектора пятнадцать равен
\begin{equation}
 \mathcal S_{\rm loop}(V_{15})=\frac17.
 \label{eq:v5-sector-selection-loop-cost}
\end{equation}
Значит, $V_{15}$ является минимумом только при заранее фиксированном
winding, но не глобальным минимумом по всем секторам.

Точный относительный вакуумный отклик также положителен:
\begin{equation}
 \Gamma_{\rm def}
 =\frac17\log\frac{m^2+a}{m^2}>0,
 \qquad a>0,quad m^2>0.
 \label{eq:v5-sector-selection-positive-response}
\end{equation}
Он измеряет цену уже имеющегося дефекта и не делает однородный сектор
неустойчивым.

Наконец, нелинейный кинк-гейт прямо установил: профиль выводится внутри
граничного сектора $Q=1$, но сами асимптотические границы были входом.
Нынешняя проверка не обнаруживает более раннего результата, который снял
бы эту зависимость.

\begin{theorem}[Допустимость и энергетическое преимущество нулевого сектора]
Текущая версия V содержит явный постоянный проекторный сектор с нулевой
кривизной и нулевым зарядом. Глобальные носители $\mathbb{RP}^3\times S^1$
и $S^4$ не индуцируют на локальной сфере свободный класс Черна один.
Положительное петлевое действие имеет абсолютное значение ноль на
тождественной петле и значение $1/7$ на минимуме сектора пятнадцать, а
относительный вакуумный отклик дефекта положителен. Поэтому существующая
топология и существующие функционалы не исключают нулевой сектор и не
порождают материю без дополнительного принципа.
\end{theorem}

\section{Отличие от уже закрытых результатов}

Результат не обнуляет класс пятнадцать. Сохраняются все строгие утверждения:
\begin{enumerate}
 \item эквивариантный точечный ёж имеет два подъёма степени $\pm1$;
 \item их хопфовы линии дают коэффициентные классы $\pm15$;
 \item внутри этого сектора совершенное операторное замыкание невозможно;
 \item минимальный положительный дефицит равен $1/7$.
\end{enumerate}
Новый no-go относится только к более сильному утверждению
\begin{equation}
 \text{``сам вакуум обязан выбрать этот сектор''}.
\end{equation}

Таким образом, топология объясняет устойчивость и устройство материи после
её появления, но пока не доказывает онтологическую неизбежность ненулевого
числа дефектов.

\section{Критерий переоткрытия}

Продолжение топологическим перебором линий, фаз и накрытий теперь является
круговым. Ветка может быть переоткрыта только одним из трёх результатов:
\begin{enumerate}
 \item полный родительский гессиан имеет отрицательную моду в нулевом
 секторе и положительный минимум в ненулевом;
 \item глобальная аномалия или индекс полной меры делает нулевой сектор
 недопустимым;
 \item исходное пространство состояний по независимой причине не содержит
 постоянной секции.
\end{enumerate}
Ни один из этих пунктов в текущем багаже не доказан.

Следующий допустимый гейт:
\path{version5_zero_sector_instability_parent_gate.tex}.

\begin{protocol}
\begin{enumerate}
 \item постоянный проекторный сектор: \textbf{существует};
 \item его локальный заряд: \textbf{ноль};
 \item $\pi_2(\mathbb{RP}^3\times S^1)$: \textbf{ноль};
 \item свободная часть $H^2(\mathbb{RP}^3\times S^1;\mathbb Z)$:
 \textbf{ноль};
 \item торсионный класс ограничивается на локальную $S^2$ нетривиально:
 \textbf{нет};
 \item $H^2(S^4;\mathbb Z)$: \textbf{ноль};
 \item нейтральность зарядов в глобально поднятой замкнутой ветви:
 \textbf{$\sum Q_i=0$};
 \item петлевое действие нулевого сектора: \textbf{ноль};
 \item петлевое действие минимума класса пятнадцать: \textbf{$1/7$};
 \item вакуумный отклик дефекта: \textbf{положителен};
 \item spin-cover bridge выбранного ежа: \textbf{сохраняется};
 \item топологический выбор ненулевого сектора: \textbf{не получен};
 \item безусловная неизбежность материи: \textbf{не доказана};
 \item следующий допустимый тест: \textbf{неустойчивость нулевого сектора}.
\end{enumerate}
\end{protocol}

Машинный сертификат сохранён в
\path{s2t_v5_equivariant_boundary_sector_selection_gate_results.json}.